\section{A Sufficient Condition for a Single Pass}
\label{sec:existence}

A link reads its context at every position $j$, but the context's final value is known only at the end of the axis.
A single pass exists when every link can be \emph{freed} of its contexts: rewritten as a function of the contexts and of reductions that read none.
This section gives the rules that free a link, the algorithm that applies them, and what follows from them.

\subsection{Freeing a Reduction from Its Contexts}
\label{sec:freeing}

Say that an operation $\varphi$ \emph{commutes} with a reduction $\oplus$ if $\varphi(a\oplus b)=\varphi(a)\oplus\varphi(b)$.
\Cref{tab:rules} lists the operations \sys pulls out of a reduction; each is an identity over $\R$ under its side condition and where all expressions are defined.
A sum distributes over addition and permits factors independent of $j$ to be extracted.
A maximum permits extraction of an added term, a factor proven non-negative, or a non-decreasing function.
An even power or an absolute value of a context-dependent term becomes the larger of two extremes.

\begin{table}[t]
\footnotesize
\caption{Operations that commute with a reduction over $j$; $a$, $c$ and $y$ do not depend on $j$. Each rule is an identity under its side condition, which \sys proves syntactically before applying it.}
\label{tab:rules}
\begin{tabular}{@{}llp{0.42\linewidth}@{}}
\toprule
\textbf{Reduction} & \textbf{Rewrite} & \textbf{Side condition}\\
\midrule
$\sum_j$ & $\sum_j (u_j+v_j)=\sum_j u_j+\sum_j v_j$ & none\\
$\sum_j$ & $\sum_j a\,u_j=a\sum_j u_j$, \ $\sum_j e^{u_j+y}=e^{y}\sum_j e^{u_j}$ & none\\
$\max_j$ & $\max_j (u_j+a)=\max_j u_j+a$ & none\\
$\max_j$ & $\max_j c\,u_j=c\max_j u_j$ & $c\ge 0$ proven\\
$\max_j$ & $\max_j h(u_j)=h(\max_j u_j)$ & $h$ non-decreasing on the relevant domain\\
$\max_j$ & $\max_j |u_j|^{2k}=\max(\max_j u_j,\ \max_j -u_j)^{2k}$ & integer $k\ge1$\\
$\min_j$ & dual to $\max_j$ & dual\\
\bottomrule
\end{tabular}
\end{table}

The separation below splits the body into $r$ contributions.
Each $g_k$ contains the data-dependent part, while $\varphi_{k,\ctx}$ carries its dependence on the context.
The context is held fixed during regrouping; the resulting summaries can therefore be accumulated before its value is known.

\begin{lemma}[Regrouping]
\label{lem:regroup}
Let $f(x,\ctx)=\bigoplus_{k=1}^{r}\varphi_{k,\ctx}(g_k(x))$, where no $g_k$ reads a context and every $\varphi_{k,\ctx}$ commutes with $\oplus$.
Then for every non-empty part $S$, $\bigoplus_{j\in S} f(x_j,\ctx)=\bigoplus_{k=1}^{r}\varphi_{k,\ctx}\bigl(\bigoplus_{j\in S} g_k(x_j)\bigr)$.
\end{lemma}

\begin{proof}
Reorder the reduction by the associativity and commutativity of $\oplus$, then pull each $\varphi_{k,\ctx}$ out.
\end{proof}

The resulting tuple can be computed without knowing the final contexts. Those contexts are used only when reconstructing the required members from the tuple.

\subsection{The Judgement}
\label{sec:judge}

A reference atom is a context-free reduction $q(S)=\bigoplus_{j\in S}g(x_j)$, identified by its reduction operation and canonical body.
Constant bodies are allowed: $|S|=\sum_{j\in S}1$ is the atom with addition and body $1$.
Before collecting atoms, every read of the reduction-axis extent as a value is normalised to this sum under the partial semantics.
Extracting a constant term from a sum can expose the same atom.
Both occurrences enter the same collection and are shared by canonical equality; no atom is added unless a required expression uses it.

\begingroup
\setlength{\intextsep}{2pt}
\reservewrapspace{16\baselineskip}
\begin{wrapfigure}{r}{0.48\textwidth}
\begin{minipage}{\linewidth}
\begin{algorithm}[H]
\small
\caption{\textsc{Judge}$(\mathit{chain}, \mathit{req})$}
\label{alg:judge}
\KwOut{reference atoms $\mathcal Q$, or failure}
Normalise extent reads in $\mathit{chain}$\;
$W\leftarrow \mathit{req}$;\ $\mathcal Q\leftarrow\emptyset$\;
\While{$W$ not empty}{
  $m\leftarrow$ pop $W$\;
  $e\leftarrow$ free contexts from body of $m$\;
  push unvisited dependencies of $e$ onto $W$\;
  \ForEach{atom $q$ in $e$}{
    \lIf{$q$ reads a context}{\Return \textsc{Stuck}$(m,q)$}
    $\mathcal Q\leftarrow\mathcal Q\cup\{q\}$ (canonical form)\;
  }
}
\Return \textsc{Fusible}$(\mathcal Q)$\;
\end{algorithm}
\end{minipage}
\end{wrapfigure}
\Cref{alg:judge} visits every required member and its dependencies, processing each member once.
It replaces each context by a symbol constant along the axis, rewrites the body with \cref{tab:rules}, and collects the remaining atoms in $\mathcal Q$.
An atom that still reads a context is \emph{stuck}: the rules could not free it.
Otherwise the collected atoms constitute the reference state tuple, with one state per distinct atom.

The construction gives a finite summary whose components join independently and whose values determine the required outputs.
Its size bounds the number of states in this construction; selection uses that bound as its default state limit.
It also identifies the context-free quantities against which a selected representation can be verified.
These roles are distinct: the existence of this construction does not establish that every representation preserving the original results fits the same limit.

\finishwrapspace
\endgroup

\begin{theorem}[Existence]
\label{thm:existence}
Suppose \textsc{Judge} returns \textsc{Fusible}$(\mathcal Q)$, with
\[
\mathcal Q=\{q_1,\dots,q_k\},\qquad
q_i(S)=\bigoplus_{j\in S}^{(i)}g_i(x_j).
\]
If the collected bodies are defined on the declared element domain, the chain has a single-pass reducer with at most $|\mathcal Q|$ states.
Its state tuple is $(q_1(S),\dots,q_k(S))$, its join is componentwise, and its epilogue reconstructs the required outputs on their source domain.
\end{theorem}

\begin{proof}[Proof sketch]
Extent normalisation preserves partial semantics, and every algebraic rewrite is an identity over $\R$ on its source domain under its proven side condition.
By \cref{lem:regroup} and induction over producer order, each required member has a reconstruction expression in the collected atoms wherever the source computation is defined.
For every atom, $q_i(S\dunion T)=q_i(S)\oplus_i q_i(T)$.
The tuple therefore joins componentwise and has leaf $(g_1(x_j),\dots,g_k(x_j))$.
The domain assumption makes these states defined on legal non-empty partitions; applying the reconstruction expression at the root gives the required outputs.
There is one coordinate per atom, establishing the bound.
\end{proof}

The judgement recognises chains whose contexts can be separated by these rules.
When a residual reduction still depends on a context, the analysis stops without producing a construction.
Further identities or stronger input relations may establish reducibility beyond this recognised class.

The construction is an algebraic reference, not necessarily the numerical form we should execute.
Separating contexts from reduction bodies changes the intermediate quantities: cancelling terms may become large independent sums, and bounded expressions may become products of unbounded factors.
Keeping the source reduction results preserves their centring and scaling, while auxiliary states provide the information needed for a join.
The next step is therefore to derive a maintained representation and prove that it retains the required information of the construction.
